% ============================================================================
%   CHAPTER 6 — CONCLUSION AND FUTURE WORK
%   Aivancity: 3-5 pages
%   UPDATED: 9 findings, external validation limitation, mode collapse
% ============================================================================

\chapter{Conclusion and Future Work}
\label{chap:conclusion}

% ----------------------------------------------------------------------------
\section{Summary of Findings}
\label{sec:conc-summary}

% [TODO — 1 page]
% Recap of the 4 pre-registered hypotheses:
%   H1: supported across all 12 configurations
%   H2: supported by 4 independent tests, main pre-registered contribution
%   H3: not reached (C-index cap at 0.65 for Meditron)
%   H4: partial (Mistral family only)

% Recap of the 9 original findings:
%   F1: GPT-4 evasion bias
%   F2: RLHF systematically degrades calibration
%   F3: Enriched-prompt paradox
%   F4: Inclusion/exclusion asymmetry
%   F5: Divergent evasion strategies across LLMs
%   F6: Post-hoc isotonic recalibration works
%   F7: Mode collapse in 5 of 12 configurations
%   F8: Mode collapse is structural, not superficial
%   F9: Global accuracy hides complete failure on discriminative cases

% ----------------------------------------------------------------------------
\section{Answer to the Research Question}
\label{sec:conc-answer}

% [TODO — 0.5 page]
% Direct answer:
% "Are the eligibility scores produced by LLM-based patient-trial matching
%  systems probabilistically reliable, and does RLHF post-training affect
%  this reliability?"
%
% ANSWER: No, they are not reliable — but the reasons are more diverse than
% initially expected. Beyond systematic miscalibration (ECE 0.12–0.63), five
% configurations exhibit mode collapse, and no medical LLM tested can handle
% truly discriminative clinical cases (0% accuracy). RLHF significantly
% degrades calibration on both tested model families. Post-hoc recalibration
% offers a practical remedy for non-degenerate models, but cannot restore
% classification quality in collapsed ones.

% ----------------------------------------------------------------------------
\section{Contributions}
\label{sec:conc-contributions}

% [TODO — 1 page]
% Nine scientific contributions:
%   1. First systematic RLHF-vs-Base calibration audit in clinical trial
%      matching (H2 supported on two independent model families)
%   2. Empirical demonstration of the enriched-prompt paradox (5 out of 6
%      models degrade under expert prompts)
%   3. Discovery of the systematic inclusion/exclusion asymmetry
%   4. Documentation of divergent evasion strategies (GPT-4 vs open-source)
%   5. Isotonic recalibration as a practical deployment recommendation,
%      with stability confirmed by cross-validation
%   6. First systematic documentation of mode collapse in open-source
%      medical LLMs, present in 5 of 12 configurations
%   7. Structural nature of mode collapse: forcing decisive labels does not
%      restore reasoning, only shifts collapse to a new default label
%   8. Complete failure of medical LLMs (Meditron, BioMistral) on truly
%      discriminative cases, hidden by global accuracy metrics
%   9. Proposal of a composite utility score combining calibration,
%      discrimination, and prediction diversity

% ----------------------------------------------------------------------------
\section{Limitations}
\label{sec:conc-limitations}

\subsection{External Validation}
\label{subsec:conc-limitations-external}

The present study relies on a single benchmark,
TrialGPT-Criterion-Annotations~\cite{jin2024trialgpt}, chosen because it is
the only public corpus offering per-criterion multi-level expert annotations
with an explicit inclusion/exclusion distinction --- properties that are
essential for fine-grained calibration analysis. External validation on
independent datasets was not carried out, for three complementary reasons.

First, existing alternatives are not directly comparable. TREC Clinical
Trials 2021/2022 and SIGIR 2016 operate at the trial level rather than the
criterion level and use only three ordinal relevance grades, which precludes
direct replication of the calibration protocol used here. Adapting the
hypotheses to these datasets would not constitute a validation but a
distinct study with different endpoints.

Second, the mode collapse phenomenon documented in
Section~\ref{sec:res-mode-collapse} appears intrinsically linked to the
availability of an evasion label (``not enough information'') in the label
space. Datasets that do not include such a label cannot be used to test
whether this behavior generalizes: they would trivially prevent it by design.

Third, prospective validation on real clinical data would require ethical
approval, hospital partnership, and integration into a live recruitment
workflow. These constraints situate external validation as a research
project in its own right, complementary to the present work rather than a
subordinate step.

External validation nevertheless remains an important direction for future
work.

\subsection{Model Scale}
\label{subsec:conc-limitations-scale}

% [TODO — 0.5 page]
% Only 7–9B models tested due to GPU constraints
% 70B medical models (Meditron-70B, Med42-v2-70B) may behave differently
% Larger models may be less prone to mode collapse

\subsection{Prompt Space}
\label{subsec:conc-limitations-prompts}

% [TODO — 0.5 page]
% Only 2 prompt variants tested
% The full prompt space is much larger (chain-of-thought, role prompting,
% etc.)
% The enriched-prompt paradox may be specific to our particular Prompt B

% ----------------------------------------------------------------------------
\section{Future Work}
\label{sec:conc-future}

\subsection{External Validation on Independent Benchmarks}
\label{subsec:conc-external}

% [TODO — 0.5 page]
% Validate findings on TREC CT (with adaptation), MIMIC-derived cohorts,
% or prospective clinical data.
% Key question: does mode collapse persist on datasets without an
% explicit NEI label?

\subsection{Multi-Agent Adversarial Debate for Calibration}
\label{subsec:conc-debate}

% [TODO — 0.5 page]
% The multi-agent debate paradigm applied to calibration.
% Cite Du et al. 2023, Khan et al. 2024.
% Hypothesis: models with distinct RLHF biases and distinct evasion
% strategies could mutually challenge each other's overconfidence and
% mode collapse via argumentative confrontation.
% To our knowledge, unexplored in clinical trial matching.

\subsection{Investigating the Origin of Mode Collapse}
\label{subsec:conc-collapse-origin}

% [TODO — 0.5 page]
% Mechanistic study: is mode collapse caused by (a) safety-oriented RLHF
% pushing models toward "safe" evasive answers, (b) domain shift between
% medical pretraining data and clinical trial criteria, or
% (c) a combination? Investigating requires model internals analysis.

\subsection{Extension to Larger and Multilingual Models}
\label{subsec:conc-scale-future}

% [TODO — 0.5 page]
% 70B models (Meditron-70B, Med42-v2-70B)
% Multilingual evaluation (MMed-Llama-3, Apollo)

% ----------------------------------------------------------------------------
\section{Closing Remarks}
\label{sec:conc-closing}

% [TODO — 0.5 page]
% Broader perspective:
% - Calibration should be a mandatory audit before clinical LLM deployment
% - The post-hoc recalibration path is practical and scalable, but only
%   for models that actually classify
% - Mode collapse is a mode of failure that current benchmarks and metrics
%   fail to detect; the community needs better evaluation protocols
% - The RLHF trade-off between usefulness and calibration deserves more
%   attention, especially in high-stakes clinical settings
